% ID: SEL-D0013-BC-E05
% ESTADO: validada. S014; tres pasadas del mismo autor, sin segundo revisor.
\textbf{Cambio de composición de la urna.}\par\nopagebreak[4]
Sean $R_i,A_i$ los colores rojo y azul en la extracción $i$. La primera bola \emph{se retira} y se sustituye por seis del color contrario, de modo que antes de la segunda extracción hay $10-1+6=15$ bolas:
\[
\begin{array}{c|rr|r}
\text{Primera extracción}&\text{Rojas después}&\text{Azules después}&\text{Total}\\\hline
R_1\ (6/10)&5&10&15\\
A_1\ (4/10)&12&3&15
\end{array}
\]
No se devuelve la bola original además de añadir seis: eso daría un modelo distinto.

\textbf{a) Segunda bola roja.}\par\nopagebreak[4]
Por probabilidad total,
\[
\boxed{P(R_2)=\frac6{10}\frac5{15}+\frac4{10}\frac{12}{15}
=\frac15+\frac8{25}=\frac{13}{25}=0.52.}
\]
\textbf{b) Primera azul, sabiendo que la segunda fue azul.}\par\nopagebreak[4]
\[
P(A_2)=1-\frac{13}{25}=\frac{12}{25},\qquad
P(A_1\cap A_2)=\frac4{10}\frac3{15}=\frac2{25}.
\]
Aplicando Bayes o la definición de condicionada,
\[
\boxed{P(A_1\mid A_2)=\frac{2/25}{12/25}=\frac16.}
\]
Las cuatro ramas conjuntas $R_1R_2,R_1A_2,A_1R_2,A_1A_2$ tienen probabilidades $0.20,0.40,0.32,0.08$ y suman uno. Las extracciones no se han tratado como independientes.
